% Pursuing Stacks (cotes 134-2 to 134-8): archive cote 134-6, site page 97 (PDF page 98), typescript page 415.
% typescript. Draft compared: ps5.tex:6474-6551.
% AI-transcribed by Claude Opus 5.5 (claude-opus-5-5, adaptive thinking, effort high), Message Batches API, batch ps3 (msgbatch_01XdFmuVVBAcgc7EFxdtCqsj), 2026-11-12; not yet checked by eye.
\scanpage{134-6/97}
\mspage{415}
most common case when this presheaf is \xout{always} a sheaf for any choice of
a,  otherwise we take the sheaf associated to it. \xout{\unsure{(NB}} In the first case
(which we may reduce to if we prefer,  by suitable choice of the site A
for given topos \hand{$\mathcal{A}$}) the functor (78) is fully faithful and moreover an
embedding,  therefore we'll identify an object a in A with the correspon-%
ding object in \hand{$\mathcal{A}$}.  Thus,  the description (76) of \uline{Hom} in terms of Hom may
be interpreted,  from this point of view,  as a formula \hand{with} X\inserted{=a} in A,
i.e. as describing the sheaf \uline{Hom}(L,F) as a functor on A$^{\circ}$. \inserted[typed]{Accordingly, the} \xout{Thus, the}
question of describing the sheaf \uline{Hom}(L,F) is reduced to the question of
describing the objects  Hom(L',F)  in \hand{$\mathcal{M}$},  for L' = $\underline{Z}^{(X)}\otimes$ L .  Thus the
main question here is to give a ``computational'' description of the objects
Hom(L,F) in \hand{$\mathcal{M}$},  for L in \hand{$\mathcal{A}$}$_{\mathrm{ab}}$  and F in \hand{$\mathcal{A}_{\mathcal{M}}$},  i.e. F and L being
sheaves on A
\[ (79) \qquad F: A^{\circ} \to \hand{$\mathcal{M}$}, \quad L : A^{\circ} \to (Ab) \quad . \]
The rule of the game here is to do so,  using just a) in case of X=a in A,
and the exactness property e).

It seems most convenient here to introduce again the \xout{category} addi-%
t\uline{i}ve envelope  Add(A) of the category A,  which we'll assume to be small
in what follows,  and the canonical \inserted[typed]{additive} functor
\[ (80) \qquad \hand{$\varepsilon_{ab}$} : Add(A) \to \hand{$\mathcal{A}$}_{ab} \quad , \]
extending the functor
\[ a \mapsto \underline{Z}^{(a)} \; : \; A \to \hand{$\mathcal{A}$}_{ab} \quad . \]
For a given F (77),  it follows from fromula (75 a)) that the \xout{corresponding}
composition
\[ \hand{$\widetilde{F}\circ\varepsilon^{\circ}_{ab}$} : Add(A)^{\circ} \xrightarrow{\hand{$\varepsilon^{\circ}_{ab}$}} \hand{$\mathcal{A}^{\circ}_{ab}$} \xrightarrow{\hand{$\widetilde{F}$}} \hand{$\mathcal{M}$} \]
is just the canonical extension \inserted{Add(F)$^{*}$} of F to Add(A)$^{\circ}$,  whose value on the gene-%
ral object
\[ x = \hand{$\bigoplus$}_{i\in I} Wh_{A}(a_{i}) \qquad \text{(I a finite indexing set)} \]
of Add(A)  (where  Wh$_{A}$(a) = $\underline{Z}^{(a)}$ as an object in Add(A)\hand{$\subset$}A\hand{$^{\wedge}$} ) is just
\[ \hand{(81} \quad \deleted{\unsure{Add(F)(x) = $\bigoplus_{i\in I}$ F(a$_{i}$)}} \quad \hand{$\mathrm{Add}(F)^{*}(x) = \prod_{i\in I} F(a_{i})$} \quad . \]
Now it is easily checked that for any object L in \hand{$\mathcal{A}$}$_{ab}$ i.e.  any sheaf
L: A$^{\circ}\to$ (Ab) ,  we have a canonical isomorphism in \hand{$\mathcal{A}$}$_{ab}$
\[ (8\unsure{1|2}) \qquad L \hand{$\xleftarrow{\sim}$} \varinjlim_{\hand{(x,u)} \text{ in } Add(A)_{/L}} \hand{$\varepsilon$}_{ab}(x) \]
(compare with the similar isomorphism on page 403).  Using (75 e),  we

\margin{97}
